% Options for packages loaded elsewhere
\PassOptionsToPackage{unicode}{hyperref}
\PassOptionsToPackage{hyphens}{url}
\documentclass[
  11pt]{ltjsarticle}
\usepackage{xcolor}
\usepackage[margin=2.5cm]{geometry}
\usepackage{amsmath,amssymb}
\setcounter{secnumdepth}{-\maxdimen} % remove section numbering
\usepackage{iftex}
\ifPDFTeX
  \usepackage[T1]{fontenc}
  \usepackage[utf8]{inputenc}
  \usepackage{textcomp} % provide euro and other symbols
\else % if luatex or xetex
  \usepackage{unicode-math} % this also loads fontspec
  \defaultfontfeatures{Scale=MatchLowercase}
  \defaultfontfeatures[\rmfamily]{Ligatures=TeX,Scale=1}
\fi
\usepackage{lmodern}
\ifPDFTeX\else
  % xetex/luatex font selection
  \ifLuaTeX
    \usepackage[hiragino-pron]{luatexja-preset}
  \fi
\fi
% Use upquote if available, for straight quotes in verbatim environments
\IfFileExists{upquote.sty}{\usepackage{upquote}}{}
\IfFileExists{microtype.sty}{% use microtype if available
  \usepackage[]{microtype}
  \UseMicrotypeSet[protrusion]{basicmath} % disable protrusion for tt fonts
}{}
\makeatletter
\@ifundefined{KOMAClassName}{% if non-KOMA class
  \IfFileExists{parskip.sty}{%
    \usepackage{parskip}
  }{% else
    \setlength{\parindent}{0pt}
    \setlength{\parskip}{6pt plus 2pt minus 1pt}}
}{% if KOMA class
  \KOMAoptions{parskip=half}}
\makeatother
\usepackage{longtable,booktabs,array}
\newcounter{none} % for unnumbered tables
\usepackage{calc} % for calculating minipage widths
% Correct order of tables after \paragraph or \subparagraph
\usepackage{etoolbox}
\makeatletter
\patchcmd\longtable{\par}{\if@noskipsec\mbox{}\fi\par}{}{}
\makeatother
% Allow footnotes in longtable head/foot
\IfFileExists{footnotehyper.sty}{\usepackage{footnotehyper}}{\usepackage{footnote}}
\makesavenoteenv{longtable}
\setlength{\emergencystretch}{3em} % prevent overfull lines
\providecommand{\tightlist}{%
  \setlength{\itemsep}{0pt}\setlength{\parskip}{0pt}}
\usepackage{bookmark}
\IfFileExists{xurl.sty}{\usepackage{xurl}}{} % add URL line breaks if available
\urlstyle{same}
\hypersetup{
  hidelinks,
  pdfcreator={LaTeX via pandoc}}

\author{}
\date{}

\begin{document}

\section{BH 熱力学プログラム 論文 9：α 自己整合方程式における 0.036
ドリフトの幾何学的起源}\label{bh-ux71b1ux529bux5b66ux30d7ux30edux30b0ux30e9ux30e0-ux8ad6ux6587-9ux3b1-ux81eaux5df1ux6574ux5408ux65b9ux7a0bux5f0fux306bux304aux3051ux308b-0.036-ux30c9ux30eaux30d5ux30c8ux306eux5e7eux4f55ux5b66ux7684ux8d77ux6e90}

\subsection{― 4 次元超球内の 2
次元面振動モードによる物理的解釈}\label{ux6b21ux5143ux8d85ux7403ux5185ux306e-2-ux6b21ux5143ux9762ux632fux52d5ux30e2ux30fcux30c9ux306bux3088ux308bux7269ux7406ux7684ux89e3ux91c8}

\textbf{著者}：木原 範昭（Noriaki Kihara） \textbf{所属}：WF System Co.,
Ltd.
\textbf{ORCID}：\href{https://orcid.org/0009-0004-6753-4020}{0009-0004-6753-4020}
\textbf{版}：v1 \textbf{日付}：2026 年 5 月 20 日
\textbf{ライセンス}：CC BY 4.0 \textbf{Concept
DOI}：\href{https://doi.org/10.5281/zenodo.20319436}{10.5281/zenodo.20319436}
\textbf{v1 Version
DOI}：\href{https://doi.org/10.5281/zenodo.20319437}{10.5281/zenodo.20319437}
\textbf{Zenodo ページ}：https://zenodo.org/records/20319437

\begin{center}\rule{0.5\linewidth}{0.5pt}\end{center}

\subsection{本稿の性格}\label{ux672cux7a3fux306eux6027ux683c}

\textbf{本稿は観察論文・解釈論文であり、第一原理証明論文ではない。}

論文 7 {[}BH7{]} の α 恒等式 \(\alpha^{-1} = 137 + (\pi^2/2)\alpha\)
における係数 \(\pi^2/2\)
の物理的起源について、標準量子論の不確定性原理と振動モードの次元解析に基づくひとつの解釈を提示する。

本稿は以下を主張しない：

\begin{itemize}
\tightlist
\item
  \(\pi^2/2\) がこの解釈以外を許容しないこと
\item
  観測値 \(\alpha^{-1} = 137.035999...\)
  の高精度残差が本稿の解釈で完全に導出されること
\item
  他の解釈の可能性を排除すること
\item
  論文 7 の主張を「証明」「強化」「拡張」すること
\end{itemize}

本稿が提示するのは、ひとつの読み方である。論文 7
の主張は本稿に依存しない。論文 7
が観察した数値的事実（\(\alpha^{-1} = 137 + (\pi^2/2)\alpha\) が 8.7 ppb
精度で成立する）は、本稿の解釈の採否によらず変わらない。

\subsection{Scope}\label{scope}

本稿の射程：

\begin{itemize}
\tightlist
\item
  \textbf{対象}：\(\alpha^{-1}(Q^2 \to 0) = 137.036\)（Thomson
  極限における α の幾何学的起源）
\item
  \textbf{方法}：標準量子論の不確定性原理＋ 4D
  空間内の振動モードの次元解析
\item
  \textbf{結論}：0.036 ドリフトは 2D
  面振動モードのゼロ点期待値として解釈可能
\end{itemize}

本稿は \textbf{以下を扱わない}：

\begin{itemize}
\tightlist
\item
  α の高エネルギー running（\(\alpha^{-1}(M_Z) \approx 127.95\)
  への流れ）
\item
  標準モデルの lepton/quark ループによる真空偏極の幾何学的再導出
\item
  結合定数の統一（α と α\_s, α\_w 等の関係）
\item
  重力との接続
\end{itemize}

これらは標準モデルの真空偏極理論で既によく記述されており、本稿の幾何学的解釈はその
\textbf{境界条件（Q² = 0 での値）}
に物理的起源を与える立場をとる。これらの高エネルギー側の幾何学的再定式化は将来課題（§8）として残す。

\subsection{要旨}\label{ux8981ux65e8}

論文 7 {[}BH7{]} において、4
次元球の単位立方体充填を用いた自己整合方程式

\[\alpha^{-1} = 137 + \frac{\pi^2}{2}\alpha\]

が、観測値 \(\alpha^{-1}(0) = 137.035999084\) と相対誤差 8.7 ppb
で一致することが報告された。本稿は、この式の構造について以下の物理的解釈を提示する：

\begin{enumerate}
\def\labelenumi{\arabic{enumi}.}
\tightlist
\item
  \textbf{α は dimensionless
  だが、物理的寄与は面積次元の量を通じて現れる}（§2）：散乱断面積
  \(\sigma_T \propto \alpha^2\) より、α の物理的役割は面積を介する。
\item
  \textbf{4D 空間の振動モードを次元別に分類すると、等方平均下で α
  に寄与しうるのは 2D 面モードのみ}（§3）：0D/1D/3D/4D
  モードは方向性を持ち平均化で消えるが、2D 面積はスカラー量として残る。
\item
  \textbf{137 個の超立方体の 2D 面が 4D
  単位球内に分布した位置位相空間の測度は \(\int_{B_4(1)} dV = \pi^2/2\)
  である}（§4）：これは数値的に \(V_4(1)\)
  と一致するが、本稿の解釈では「2D 面の位置自由度の積分」として読む。
\item
  \textbf{W7 自己整合方程式 \(\alpha^{-1} = 137 + (\pi^2/2)\alpha\)
  は、137 個の超立方体の 2D
  面モードのゼロ点振動による自己整合的補正として解釈できる}（§5）。
\item
  \textbf{この構造は Wilson
  格子ゲージ理論のプラケット作用と同一}（§6）：論文 8 {[}BH8{]}
  で示された同型性が幾何学的レベルで具体化される。
\item
  \textbf{α の観測値には、2D
  モード分布の幅に由来する理論的下限が存在}（§7）。
\end{enumerate}

本稿の解釈は標準量子論の枠組み内で完結し、新たな仮定（離散時空、複素整数格子等）を必要としない。

\begin{center}\rule{0.5\linewidth}{0.5pt}\end{center}

\subsection{§1 導入}\label{ux5c0eux5165}

論文 7 {[}BH7{]} は、4 次元球 \(B_4(R=3)\) に内接する 137 個の単位 4
次元立方体（テセラクト）と、4 次元単位球の体積 \(V_4(1) = \pi^2/2\)
を用いて、α の自己整合方程式

\[\alpha^{-1} = N(1) + V_4(1)\alpha = 137 + \frac{\pi^2}{2}\alpha\]

を提示した。この方程式は観測値 \(\alpha^{-1}(0) = 137.035999084(21)\) と
8.7 ppb の精度で一致する。

この自己整合方程式の構造には、以下の未解決問題があった：

\textbf{問題 1}：\(\pi^2/2\) がなぜ α の係数として現れるか - 数値的には
\(V_4(1)\) と一致するが、なぜ 4D
単位球の体積が登場するのか物理的に不明確 -
単なる数学的同一性なのか、物理的意味があるのか

\textbf{問題 2}：0.036 ドリフトの幾何学的起源 - 137
という整数が出る幾何学的構造（W6 論文）は明確 -
しかし「0.036」の起源は自己整合方程式の閉包としてのみ説明されていた

\textbf{問題 3}：W8 {[}BH8{]} で示された Wilson
格子ゲージ理論との同型性の物理的内容 - 構造的対応は確立 -
しかし対応の物理的内容（なぜそうなるか）は未解明

本稿は、これらの問題に対し以下の解釈を提示する：

\begin{quote}
\textbf{α
は物理的に面積次元の量（\(\sigma \propto \alpha^2\)）であり、4D
空間の振動モードのうち α に寄与しうるのは 2D 面モードのみ。137
個の超立方体の 2D 面振動モードの位置位相空間測度が \(\pi^2/2\)
を与え、これが W7 自己整合方程式の係数構造を物理的に説明する。}
\end{quote}

この解釈は標準量子論の不確定性原理を出発点とし、新たな仮定を必要としない。

\begin{center}\rule{0.5\linewidth}{0.5pt}\end{center}

\subsection{§2 α
の次元解析と面積次元性}\label{ux3b1-ux306eux6b21ux5143ux89e3ux6790ux3068ux9762ux7a4dux6b21ux5143ux6027}

\subsubsection{2.1 α 自身は
dimensionless}\label{ux3b1-ux81eaux8eabux306f-dimensionless}

微細構造定数 α は SI 単位系で

\[\alpha = \frac{e^2}{4\pi\epsilon_0 \hbar c} \approx 7.297 \times 10^{-3}\]

として定義され、dimensionless である。

\subsubsection{2.2 物理的寄与は α²
を通じて現れる}\label{ux7269ux7406ux7684ux5bc4ux4e0eux306f-ux3b1uxb2-ux3092ux901aux3058ux3066ux73feux308cux308b}

しかし α が物理的に関与する量は、多くの場合 \(\alpha^2\) で現れる：

\textbf{Thomson 散乱断面積}：
\[\sigma_T = \frac{8\pi}{3} r_e^2 = \frac{8\pi}{3} \left(\alpha \cdot \frac{\hbar}{mc}\right)^2 \propto \alpha^2 \cdot [\text{長さ}]^2\]

\textbf{Rutherford 散乱微分断面積}：
\[\frac{d\sigma}{d\Omega} \propto \alpha^2 \cdot [\text{長さ}]^2\]

\textbf{Bohr 半径}：
\[a_0 = \frac{\hbar}{\alpha mc} \propto \alpha^{-1} \cdot [\text{長さ}]\]

\textbf{古典電子半径}：
\[r_e = \alpha \cdot \frac{\hbar}{mc} \propto \alpha \cdot [\text{長さ}]\]

すなわち、α は \textbf{断面積（面積次元）} と直結する量である。

\subsubsection{2.3 4D
空間内の幾何学的構造との対応}\label{d-ux7a7aux9593ux5185ux306eux5e7eux4f55ux5b66ux7684ux69cbux9020ux3068ux306eux5bfeux5fdc}

4D 空間内で α が結合する自然な幾何学的構造を考えると：

\begin{itemize}
\tightlist
\item
  σ（断面積）∝ α² ∝ 面積
\item
  → α は \textbf{2 次元的な構造}に結合する量
\end{itemize}

これは Wilson 格子ゲージ理論において、ゲージ結合定数がプラケット（2D
面、2-form）の作用係数として現れることと整合する。

\subsubsection{2.4
本稿の出発点}\label{ux672cux7a3fux306eux51faux767aux70b9}

以上より：

\begin{quote}
\textbf{α への寄与を 4D 空間の振動モードから考える場合、α
の面積次元性により、寄与しうるのは「2
次元的な面積を持つ振動モード」に限定される。}
\end{quote}

これが本稿の核心的観察である。

\begin{center}\rule{0.5\linewidth}{0.5pt}\end{center}

\subsection{§3 4D
空間内の振動モードの次元別分類}\label{d-ux7a7aux9593ux5185ux306eux632fux52d5ux30e2ux30fcux30c9ux306eux6b21ux5143ux5225ux5206ux985e}

\subsubsection{3.1 137 個の超立方体と外接
4-球面の振動モード}\label{ux500bux306eux8d85ux7acbux65b9ux4f53ux3068ux5916ux63a5-4-ux7403ux9762ux306eux632fux52d5ux30e2ux30fcux30c9}

論文 7 で扱う系： - 半径 \(R = 3\) の 4 次元球 \(B_4(3)\) -
その内部に充填された 137 個の単位 4 次元立方体（テセラクト）

この系には、以下の振動モードが存在する：

\textbf{テセラクト 1 個あたりの内部構造}：

{\def\LTcaptype{none} % do not increment counter
\begin{longtable}[]{@{}lll@{}}
\toprule\noalign{}
構造要素 & 個数 & 振動モード数（独立な） \\
\midrule\noalign{}
\endhead
\bottomrule\noalign{}
\endlastfoot
頂点（0-cell） & 16 & 16 × 4 = 64 方向 \\
辺（1-cell） & 32 & 伸縮 32、横振動 等 \\
面（2-cell、正方形） & 24 & 膨らみ、捻じれ \\
立体（3-cell、立方体） & 8 & 体積振動 \\
4-胞（4-cell、テセラクト全体） & 1 & 全体呼吸モード \\
\end{longtable}
}

\textbf{137 個全体}： - 各構造要素の振動が 137 倍 -
隣接立方体間の相互作用モード - 全体としてのフォノン的スペクトル

\textbf{外接 4-球面 \(S^3\)（半径 3）}：

{\def\LTcaptype{none} % do not increment counter
\begin{longtable}[]{@{}
  >{\raggedright\arraybackslash}p{(\linewidth - 2\tabcolsep) * \real{0.4324}}
  >{\raggedright\arraybackslash}p{(\linewidth - 2\tabcolsep) * \real{0.5676}}@{}}
\toprule\noalign{}
\begin{minipage}[b]{\linewidth}\raggedright
モード
\end{minipage} & \begin{minipage}[b]{\linewidth}\raggedright
説明
\end{minipage} \\
\midrule\noalign{}
\endhead
\bottomrule\noalign{}
\endlastfoot
動径方向呼吸 & R = 3 の振動 \\
球面調和 \(Y_{\ell m n}\) on \(S^3\) &
角度方向の振動（\(\ell = 0, 1, 2, \ldots\)） \\
4 軸方向の歪み & 球 → 楕円体変形 \\
\end{longtable}
}

\subsubsection{3.2
等方平均下での生存条件}\label{ux7b49ux65b9ux5e73ux5747ux4e0bux3067ux306eux751fux5b58ux6761ux4ef6}

これらのモードのうち、α への寄与を考える際に重要なのは
\textbf{等方平均下での生存条件} である。

α は dimensionless かつスカラーであり、SO(4) 対称性下で不変。したがって
α への寄与は、\textbf{等方平均（4
軸方向の平均化）下でゼロにならないモード}に限定される。

{\def\LTcaptype{none} % do not increment counter
\begin{longtable}[]{@{}
  >{\raggedright\arraybackslash}p{(\linewidth - 4\tabcolsep) * \real{0.1301}}
  >{\raggedright\arraybackslash}p{(\linewidth - 4\tabcolsep) * \real{0.5285}}
  >{\raggedright\arraybackslash}p{(\linewidth - 4\tabcolsep) * \real{0.3415}}@{}}
\toprule\noalign{}
\begin{minipage}[b]{\linewidth}\raggedright
モード次元
\end{minipage} & \begin{minipage}[b]{\linewidth}\raggedright
数学的性質
\end{minipage} & \begin{minipage}[b]{\linewidth}\raggedright
等方平均の結果
\end{minipage} \\
\midrule\noalign{}
\endhead
\bottomrule\noalign{}
\endlastfoot
0D（点） & ベクトル位置の揺らぎ & \textbf{ゼロ}（並進対称性） \\
1D（線分） & 方向を持つ &
\textbf{ゼロ}（線対称性、向きが反転すると相殺） \\
\textbf{2D（面）} & \textbf{面積（向きを持たないスカラー量）} &
\textbf{生き残る} \\
3D（体積、有向） & 体積要素 \(dV = dx \wedge dy \wedge dz\)
は擬スカラー（向きを持つ） & \textbf{ゼロ}（向き反転で符号反転） \\
4D（4-体積） & 4-体積要素は擬スカラー & \textbf{ゼロ} \\
\end{longtable}
}

\textbf{鍵となる事実}：\textbf{面積はスカラー量}であり、面の向きを反転しても面積は不変。一方、長さ・体積（有向な体積要素）は向きを反転すると符号が変わる。

\subsubsection{3.3 選別原理}\label{ux9078ux5225ux539fux7406}

以上より：

\begin{quote}
\textbf{4D 空間内の振動モードのうち、等方平均下で α
に有意な寄与を与えうるのは 2D 面モードのみである。}
\end{quote}

これは α の面積次元性（§2）と整合する。

\subsubsection{3.4 Wilson
格子ゲージ理論との整合}\label{wilson-ux683cux5b50ux30b2ux30fcux30b8ux7406ux8ad6ux3068ux306eux6574ux5408}

この結論は Wilson 格子ゲージ理論における \textbf{2-form
ゲージ場（プラケット作用）} の構造と一致する：

\begin{itemize}
\tightlist
\item
  0-form（スカラー場）：ゲージ不変な点的構造
\item
  1-form（ベクトルポテンシャル \(A_\mu\)）：ゲージ依存
\item
  \textbf{2-form（場の強度 \(F_{\mu\nu}\)）：ゲージ不変、プラケット作用
  \(S = \beta \sum_p \text{Re}\,\text{tr}(U_p)\)}
\item
  3-form, 4-form：高次の構造
\end{itemize}

α が「面に結合する」ことは、Wilson 形式での標準的事実である。

\begin{center}\rule{0.5\linewidth}{0.5pt}\end{center}

\subsection{§4 137 個の超立方体の 2D
面モード位置位相空間}\label{ux500bux306eux8d85ux7acbux65b9ux4f53ux306e-2d-ux9762ux30e2ux30fcux30c9ux4f4dux7f6eux4f4dux76f8ux7a7aux9593}

\subsubsection{4.1 2D
面モードの位置自由度}\label{d-ux9762ux30e2ux30fcux30c9ux306eux4f4dux7f6eux81eaux7531ux5ea6}

2D 面振動モードを完全に指定するには、以下が必要：

\begin{itemize}
\tightlist
\item
  \textbf{位置}：4 次元球 \(B_4(R=3)\) 内のどこに 2D 面があるか
\item
  \textbf{方向}：4 次元空間内の 6 個の
  2-平面方向（\(\binom{4}{2} = 6\)）のうちどれか
\item
  \textbf{振幅}：振動の励起度
\end{itemize}

\subsubsection{4.2
位置自由度の積分}\label{ux4f4dux7f6eux81eaux7531ux5ea6ux306eux7a4dux5206}

固定された方向と振幅に対して、2D 面の位置自由度を 4
次元単位球（スケール正規化済み）内で積分すると：

\[\text{2D 面の位置位相空間} = \int_{B_4(1)} dV = \frac{\pi^2}{2}\]

これは数値的に \textbf{\(V_4(1) = \pi^2/2\)} と一致する。

\subsubsection{4.3
本稿の解釈：物理的意味の読み替え}\label{ux672cux7a3fux306eux89e3ux91c8ux7269ux7406ux7684ux610fux5473ux306eux8aadux307fux66ffux3048}

{\def\LTcaptype{none} % do not increment counter
\begin{longtable}[]{@{}
  >{\raggedright\arraybackslash}p{(\linewidth - 2\tabcolsep) * \real{0.2297}}
  >{\raggedright\arraybackslash}p{(\linewidth - 2\tabcolsep) * \real{0.7703}}@{}}
\toprule\noalign{}
\begin{minipage}[b]{\linewidth}\raggedright
解釈
\end{minipage} & \begin{minipage}[b]{\linewidth}\raggedright
\(\pi^2/2\) の物理的意味
\end{minipage} \\
\midrule\noalign{}
\endhead
\bottomrule\noalign{}
\endlastfoot
旧解釈（W7 論文） & 4 次元単位球の体積 \(V_4(1)\) \\
\textbf{本稿の解釈} & \textbf{2D 面モードを 4D
単位球内に配置できる位置自由度の測度} \\
\end{longtable}
}

両者は
\textbf{数値的に同一だが、物理的意味は異なる}。本稿の解釈の下では、\(\pi^2/2\)
は α の係数として現れる必然性を持つ：α が 2D
面モードに結合する以上、その位置位相空間体積が結合定数の幾何学的係数として現れる。

\subsubsection{4.4
方向と振幅の寄与}\label{ux65b9ux5411ux3068ux632fux5e45ux306eux5bc4ux4e0e}

2D 面の方向（6 個の
2-平面）の寄与は、組み合わせ的因子として別途現れうる。本稿では位置積分のみを主要因子として扱い、方向因子は数値的に効くとしても
\(O(1)\) の補正と見なす。

振幅は α 自身（励起確率振幅）として現れる：

\[\text{2D 面モードの集団的寄与} = \underbrace{\frac{\pi^2}{2}}_{\text{位置位相空間}} \times \underbrace{\alpha}_{\text{各面の振幅}} = \frac{\pi^2}{2} \alpha\]

これが W7 自己整合方程式における α の係数項を物理的に説明する。

\begin{center}\rule{0.5\linewidth}{0.5pt}\end{center}

\subsection{§5 W7
自己整合方程式の物理的解釈}\label{w7-ux81eaux5df1ux6574ux5408ux65b9ux7a0bux5f0fux306eux7269ux7406ux7684ux89e3ux91c8}

\subsubsection{5.1
方程式の構造}\label{ux65b9ux7a0bux5f0fux306eux69cbux9020}

論文 7 {[}BH7{]} の自己整合方程式：

\[\alpha^{-1} = 137 + \frac{\pi^2}{2}\alpha\]

\subsubsection{5.2
本稿の解釈による各項}\label{ux672cux7a3fux306eux89e3ux91c8ux306bux3088ux308bux5404ux9805}

{\def\LTcaptype{none} % do not increment counter
\begin{longtable}[]{@{}
  >{\raggedright\arraybackslash}p{(\linewidth - 4\tabcolsep) * \real{0.1304}}
  >{\raggedright\arraybackslash}p{(\linewidth - 4\tabcolsep) * \real{0.3354}}
  >{\raggedright\arraybackslash}p{(\linewidth - 4\tabcolsep) * \real{0.5342}}@{}}
\toprule\noalign{}
\begin{minipage}[b]{\linewidth}\raggedright
項
\end{minipage} & \begin{minipage}[b]{\linewidth}\raggedright
数学的内容
\end{minipage} & \begin{minipage}[b]{\linewidth}\raggedright
物理的解釈（本稿）
\end{minipage} \\
\midrule\noalign{}
\endhead
\bottomrule\noalign{}
\endlastfoot
\textbf{137} & 4D 球 \(B_4(3)\) に内接する単位立方体充填数 &
\textbf{全振動モードが基底状態（最低エネルギー）にある場合の寄与}：4D ℤ⁴
格子の整数論的事実 \\
\textbf{\(\pi^2/2\)} & \(V_4(1)\)、または本稿の解釈では 2D
面モード位置位相空間 & \textbf{137 個の超立方体の 2D 面モードを 4D
単位球内に配置できる位置自由度} \\
\textbf{α}（係数項の） & 結合定数 & \textbf{各 2D 面モードの励起振幅} \\
\textbf{\((\pi^2/2)\alpha\)} & 0.036（数値的に） & \textbf{2D
面モードのゼロ点振動の集団的寄与}（標準量子論の不確定性原理から） \\
\end{longtable}
}

\subsubsection{5.3
標準量子論の不確定性原理との接続}\label{ux6a19ux6e96ux91cfux5b50ux8ad6ux306eux4e0dux78baux5b9aux6027ux539fux7406ux3068ux306eux63a5ux7d9a}

各 2D 面モードは、調和振動子近似下で以下のゼロ点エネルギーを持つ：

\[E_{\text{zero-point}}^{(i)} = \frac{1}{2}\hbar\omega_i\]

ここで \(\omega_i\) は i 番目の 2D 面モードの固有振動数。

全 2D 面モードの集団的寄与は：

\[\sum_i \frac{1}{2}\hbar\omega_i = (\text{位置位相空間測度}) \times (\text{振幅密度}) = \frac{\pi^2}{2} \cdot \alpha\]

ここで右辺の α は、各モードの平均励起振幅が自己整合的に α
自身に等しくなる、という自己整合条件を表す。

\subsubsection{5.4 自己整合性}\label{ux81eaux5df1ux6574ux5408ux6027}

α が 2D 面モードの振幅であり、かつ 2D 面モードのゼロ点寄与が α⁻¹
に補正を与えるという構造から、自己整合方程式

\[\alpha^{-1} = 137 + \frac{\pi^2}{2}\alpha \quad \Leftrightarrow \quad \frac{\pi^2}{2}\alpha^2 + 137\alpha - 1 = 0\]

が自然に導かれる。

\subsubsection{5.5 数値的一致}\label{ux6570ux5024ux7684ux4e00ux81f4}

この方程式の正根：

\[\alpha^{-1} = 137.036010988\ldots\]

CODATA 2018 値 \(\alpha^{-1}(0) = 137.035999084(21)\) との相対誤差は
\textbf{8.7 ppb}。これは論文 7
で既に観察された事実であり、本稿の解釈はこの数値的一致を説明する
\textbf{物理的フレームワーク}を提供する。

\begin{center}\rule{0.5\linewidth}{0.5pt}\end{center}

\subsection{§6 Wilson
格子ゲージ理論との同型性の強化}\label{wilson-ux683cux5b50ux30b2ux30fcux30b8ux7406ux8ad6ux3068ux306eux540cux578bux6027ux306eux5f37ux5316}

論文 8 {[}BH8{]} は、W7 自己整合方程式と Wilson 格子ゲージ理論の間に
Schläfli 双対と B\_4 同変性に基づく構造的対応があることを示した。本稿の
2D 面モード解釈により、この対応の \textbf{物理的内容}が明確化される。

\subsubsection{6.1 対応関係}\label{ux5bfeux5fdcux95a2ux4fc2}

{\def\LTcaptype{none} % do not increment counter
\begin{longtable}[]{@{}
  >{\raggedright\arraybackslash}p{(\linewidth - 2\tabcolsep) * \real{0.5978}}
  >{\raggedright\arraybackslash}p{(\linewidth - 2\tabcolsep) * \real{0.4022}}@{}}
\toprule\noalign{}
\begin{minipage}[b]{\linewidth}\raggedright
Wilson 格子ゲージ理論
\end{minipage} & \begin{minipage}[b]{\linewidth}\raggedright
W7 自己整合（本稿の解釈）
\end{minipage} \\
\midrule\noalign{}
\endhead
\bottomrule\noalign{}
\endlastfoot
結合定数 \(\beta = 1/g^2\) & \(\alpha^{-1}\) \\
プラケット（2-form 構造） & 2D 面振動モード \\
プラケット数 \(\sum_p\) & \(\pi^2/2\)（2D 面の位置位相空間測度） \\
プラケット作用 \(\beta \sum_p \text{Re}\,\text{tr}(U_p)\) &
\((\pi^2/2)\alpha\) \\
Cell（4D の Hard core） & 137（単位立方体充填数） \\
全作用 & \(\alpha^{-1} = 137 + (\pi^2/2)\alpha\) \\
\end{longtable}
}

\subsubsection{6.2
物理的内容の明確化}\label{ux7269ux7406ux7684ux5185ux5bb9ux306eux660eux78baux5316}

論文 8 の同型性は
\textbf{構造的対応}として確立されていたが、本稿により：

\begin{quote}
\textbf{W7 自己整合方程式は、Wilson
格子ゲージ理論のプラケット作用構造の、4D 球幾何上での具体化である。}
\end{quote}

として \textbf{物理的に位置付け} られる。

\subsubsection{6.3 含意}\label{ux542bux610f}

これにより： - W7 と W8
の対応が「数学的同型」から「物理的同型」へと格上げ - α が「2D
面に結合する」という Wilson 形式の自然な解釈と一致 - 標準 QFT
の真空偏極を \textbf{幾何学的に再構成}するロードマップが見える

\begin{center}\rule{0.5\linewidth}{0.5pt}\end{center}

\subsection{§7 α
の観測値の理論的下限幅}\label{ux3b1-ux306eux89b3ux6e2cux5024ux306eux7406ux8ad6ux7684ux4e0bux9650ux5e45}

\subsubsection{7.1
ゼロ点振動の分布幅}\label{ux30bcux30edux70b9ux632fux52d5ux306eux5206ux5e03ux5e45}

各 2D
面モードのゼロ点振動は、標準量子論の不確定性原理から有限の分布幅を持つ：

\[\Delta n_i^{\text{zero}} = \frac{1}{2}\]

（光子数演算子のゼロ点ゆらぎ）

\subsubsection{7.2 α への伝播}\label{ux3b1-ux3078ux306eux4f1dux64ad}

全 2D 面モードの集団的寄与に対するゆらぎは：

\[(\Delta \alpha^{-1})^2 \approx \sum_i \left(\frac{\partial \alpha^{-1}}{\partial n_i}\right)^2 (\Delta n_i)^2\]

具体的な計算は本稿の射程外（モード固有振動数の決定が必要）だが、構造的には：

\begin{quote}
\textbf{α の観測値には、2D
面モード分布の幅に由来する理論的下限が存在する。}
\end{quote}

\subsubsection{7.3 CODATA
実測値との対応}\label{codata-ux5b9fux6e2cux5024ux3068ux306eux5bfeux5fdc}

\[\alpha^{-1}(0) = 137.035999084(21)\]

CODATA の不確定性 \((21) \approx 2.1 \times 10^{-8}\)
は実験精度であり、本稿の理論的下限がこれよりさらに小さいかどうかは、モード振動数のさらなる検討を要する。

\subsubsection{7.4 含意}\label{ux542bux610f-1}

本稿の解釈により： - α は「鋭い点」ではなく「分布の中心値」 - 観測される
8.7 ppb のずれ（W7）は、モード寄与の高次補正の余地 -
将来の精密測定で本稿の理論的下限が観測可能になる可能性

\begin{center}\rule{0.5\linewidth}{0.5pt}\end{center}

\subsection{§8
残された問題：高エネルギーへの拡張}\label{ux6b8bux3055ux308cux305fux554fux984cux9ad8ux30a8ux30cdux30ebux30aeux30fcux3078ux306eux62e1ux5f35}

本稿の解釈は \textbf{Thomson 極限 \(\alpha^{-1}(Q^2 \to 0) = 137.036\)
のみを扱う}。

\subsubsection{8.1 標準 QED における α の
running}\label{ux6a19ux6e96-qed-ux306bux304aux3051ux308b-ux3b1-ux306e-running}

標準モデルでは、α はエネルギースケール \(Q^2\) に依存して走る：

\begin{itemize}
\tightlist
\item
  \(\alpha^{-1}(0) = 137.036\)（Thomson 極限）
\item
  \(\alpha^{-1}(M_Z^2) \approx 127.95\)（Z ボソン質量スケール）
\item
  差分 \(\approx 9.08\) 単位は、lepton/quark loop
  による真空偏極で説明される
\end{itemize}

この running は標準 QED の精密測定で
\textbf{十分よく記述}されており、本稿は競合しない。

\subsubsection{8.2 本稿の立場}\label{ux672cux7a3fux306eux7acbux5834}

本稿の幾何学的解釈は、α の \textbf{境界条件 \(Q^2 = 0\)
での値}に物理的起源を与える。 running は標準 QED に委譲する。

両者は階層的に補完：

{\def\LTcaptype{none} % do not increment counter
\begin{longtable}[]{@{}lll@{}}
\toprule\noalign{}
階層 & 内容 & 担当 \\
\midrule\noalign{}
\endhead
\bottomrule\noalign{}
\endlastfoot
1 & 137（整数、幾何不変） & 論文 6 {[}BH6{]}：充填の整数論 \\
2 & \((\pi^2/2)\alpha \approx 0.036\) & \textbf{本論文：2D 面モード} \\
3 & \(\Delta\alpha_{SM}(Q^2)\) & 標準 QED 真空偏極 \\
\end{longtable}
}

\subsubsection{8.3 将来課題}\label{ux5c06ux6765ux8ab2ux984c}

階層 3（標準 QED の真空偏極）と本稿の幾何学的解釈との関係は、以下の open
questions を残す：

\begin{enumerate}
\def\labelenumi{\arabic{enumi}.}
\tightlist
\item
  \textbf{高エネルギーで活性化される追加モード}：3D, 4D
  振動モードが高エネルギーで α への寄与を始める可能性
\item
  \textbf{SM ループの幾何学的対応物}：electron loop, muon loop, quark
  loop の各々が、特定の幾何学的モードに対応するか
\item
  \textbf{\(\alpha^{-1}(M_Z) \approx 128\)
  の幾何学的導出}：本稿の框架を高エネルギーに拡張できるか
\item
  \textbf{GUT スケールでの α の挙動}：grand unification scale
  での結合定数の幾何学的予言
\end{enumerate}

これらの幾何学的再定式化は将来研究の課題である。\textbf{本稿の主張（Thomson
極限の幾何学的起源）は、running の機構が何であれ独立に成立する}。

\begin{center}\rule{0.5\linewidth}{0.5pt}\end{center}

\subsection{§9 結論}\label{ux7d50ux8ad6}

本稿は、論文 7 {[}BH7{]} の α 自己整合方程式
\(\alpha^{-1} = 137 + (\pi^2/2)\alpha\)
の物理的解釈として、以下を提示した：

\begin{enumerate}
\def\labelenumi{\arabic{enumi}.}
\tightlist
\item
  \textbf{α は dimensionless
  だが物理的に面積次元の量}（\(\sigma \propto \alpha^2\)）
\item
  \textbf{4D 空間の振動モードのうち、等方平均下で α に寄与しうるのは 2D
  面モードのみ}
\item
  \textbf{137 個の超立方体の 2D 面が 4D
  単位球内に分布した位置位相空間の測度 = \(\pi^2/2\)}
\item
  \textbf{W7 自己整合方程式は 2D
  面モードのゼロ点振動の自己整合的寄与として解釈可能}
\item
  \textbf{この構造は Wilson 格子ゲージ理論のプラケット作用と同一}（W8
  同型性の物理的内容）
\item
  \textbf{観測される α には 2D モード分布幅に由来する理論的下限が存在}
\end{enumerate}

本稿の解釈は標準量子論の不確定性原理の枠内で完結し、新たな仮定（離散時空、追加次元等）を必要としない。

\textbf{本稿の射程}は α の Thomson 極限のみ。高エネルギーへの running
は標準 QED に委譲し、その幾何学的再定式化は将来課題（§8）として残す。

\subsubsection{本稿の position}\label{ux672cux7a3fux306e-position}

論文 7 {[}BH7{]}、論文 8 {[}BH8{]}、本論文の三者の関係：

{\def\LTcaptype{none} % do not increment counter
\begin{longtable}[]{@{}
  >{\raggedright\arraybackslash}p{(\linewidth - 4\tabcolsep) * \real{0.2564}}
  >{\raggedright\arraybackslash}p{(\linewidth - 4\tabcolsep) * \real{0.5641}}
  >{\raggedright\arraybackslash}p{(\linewidth - 4\tabcolsep) * \real{0.1795}}@{}}
\toprule\noalign{}
\begin{minipage}[b]{\linewidth}\raggedright
論文
\end{minipage} & \begin{minipage}[b]{\linewidth}\raggedright
内容
\end{minipage} & \begin{minipage}[b]{\linewidth}\raggedright
担当領域
\end{minipage} \\
\midrule\noalign{}
\endhead
\bottomrule\noalign{}
\endlastfoot
BH7（α 恒等式） & \(\alpha^{-1} = 137 + (\pi^2/2)\alpha\) の発見 &
代数的観察 \\
BH8（Wilson 同型性） & 構造的対応の証明 & 数学的対応 \\
\textbf{本論文（BH9）} & \textbf{2D 面モードによる物理的解釈} &
\textbf{物理的内容} \\
\end{longtable}
}

この 3 部作により、α の幾何学的起源は \textbf{観察 → 構造 →
物理的内容}の 3 層で位置付けられる。

\begin{center}\rule{0.5\linewidth}{0.5pt}\end{center}

\subsection{参考文献}\label{ux53c2ux8003ux6587ux732e}

{[}BH6{]} Kihara, N. (2026). \emph{Synthesis and the 4+1 to 3+1
Reduction}. Zenodo. Concept DOI: 10.5281/zenodo.19837597.

{[}BH7{]} Kihara, N. (2026). \emph{A Geometric Identity for the
Fine-Structure Constant: From the 4D Unit Ball Volume and its
Cube-Packing Deficit}. Zenodo. Concept DOI: 10.5281/zenodo.19869266.

{[}BH8{]} Kihara, N. (2026). \emph{Chain Complex Structure on the 4D
Hypercubic Lattice: Structural Correspondence between Kihara
Cube-Packing and Wilson Lattice Gauge Theory}. Zenodo. Concept DOI:
10.5281/zenodo.19880467.

{[}BH7-Supp{]} Kihara, N. (2026). \emph{論文 7 補講：α
恒等式の二次補正項に関する幾何学的観察}. Zenodo. Concept DOI:
10.5281/zenodo.19933729.

{[}Wilson1974{]} Wilson, K. G. (1974). \emph{Confinement of quarks}.
Phys. Rev.~D \textbf{10}, 2445.

{[}CODATA2018{]} Tiesinga, E. et al.~(2021). \emph{CODATA Recommended
Values of the Fundamental Physical Constants: 2018}. Rev.~Mod. Phys.
\textbf{93}, 025010.

\begin{center}\rule{0.5\linewidth}{0.5pt}\end{center}

\subsection{改訂履歴}\label{ux6539ux8a02ux5c65ux6b74}

\begin{itemize}
\tightlist
\item
  \textbf{v1（2026 年 5 月 20 日）}：初版
\end{itemize}

\end{document}
